(i) Since $\omega \circ \alpha_{g} = \omega$ and $\alpha_{g}(O_{i}) = O_{gi}$, $C_{gi,gj} = C_{ij}$ for all $g \in G$. Transitivity on the labels gives $C_{ii} = v$ for all $i$. (For $N \ge 3$ transitivity on unordered pairs implies it: if the labels split into two nonempty invariant sets, the pairs inside one set and the pairs across the two are different orbits, and one set has two labels or the other does.) $C$ is symmetric, so $C_{ij}$ is a function of the unordered pair $\{i, j\}$, constant on the single orbit: $C_{ij} = c$ for $i \ne j$. Then $C = (v - c)1 + cU$; $U$ has eigenvalue $N$ on the constants and $0$ on the vectors of zero sum, which gives the two eigenvalues, and $C \ge 0$ (a covariance) is $v - c \ge 0$ and $v + (N - 1)c \ge 0$. The check computes the orbits by closure for the symmetric group on four labels, the affine groups of the fields of five and seven elements (doubly transitive), the group $x \mapsto a^{2}x + b$ of seven elements (transitive on unordered pairs with two orbits on ordered pairs), and the cyclic and dihedral groups of six (three orbits each), and verifies the eigenvalues and the positivity ends.

(ii) $\phi(f) = \sum_{i}f_{i}(O_{i} - \omega(O_{i}))$, so $\mathrm{Var}_{\omega}(\phi(f)) = \sum_{ij}f_{i}C_{ij}f_{j} = (v - c)\sum_{i}f_{i}^{2} + c(\sum_{i}f_{i})^{2}$ by (i). If the profiles tile the slice, $\sum_{i}g_{i} = 1$, then $\sum_{i}f_{i} = \int f\sum_{i}g_{i} = \int f$, which vanishes for a cosine or sine mode of nonzero wave vector; the remaining term $(v - c)\sum_{i}f_{i}^{2}$ depends on $k$ only through the transforms of the cells, which tend to their volumes for $k$ below the inverse cell size. On a periodic lattice a cosine or sine mode of nonzero dual wave vector sums to zero over the cells, and $\sum_{i}f_{i}^{2}$ is the number of cells times the mean square of the mode, so the variance per unit volume and unit mean square is $(v - c)V_{c}$, independent of the wave vector. A spectrum independent of $k$ has $\Delta^{2} \propto k^{3}$, the limiting form of the stored bound of Theorem \claimref{synthesis.finite-range-descent-deposit-cannot-seed-cmb}(i) for a covariance supported within one cell. For point cells at arbitrary points, $f_{i} = V_{c}\cos(k\cdot x_{i})$ or $V_{c}\sin(k\cdot x_{i})$, and adding the two variances gives $V_{c}^{2}[(v - c)N + c\,((\sum_{i}\cos k\cdot x_{i})^{2} + (\sum_{i}\sin k\cdot x_{i})^{2})]$, which is the boundary statement; the second bracket vanishes at nonzero $k$ exactly on the scales where the points have no structure factor, as on a lattice away from its dual. For two cells at $0$ and $L/4$ the cosines are $(1, 0)$ and $(1, -1)$, and $f^{T}Cf$ is $v$ and $2(v - c)$. The check verifies the equal variances on the three cosine modes of a ring of six cells (the cosines are fixed by their recurrence, with rational values) and on the seven Walsh modes of eight cells, for several $(v, c)$. A further check function takes a ring of six sites: point cells on the whole ring and on a sublattice (no $c$ term except on the dual lattice), cells that tile the ring with equal and with unequal sizes (no $c$ term at nonzero $k$, and $(v - c)$ times the cells' form factor), point cells at the sites $0, 1, 3$ (the $c$ term is $c$ times $N S_{\mathrm{pl}} = 1, 3, 1$ on the three modes, so the first two modes agree only for $c = 0$), and the two cells a quarter circle apart.

(iii) $A$ commutes with the permutation matrices of $G$ because $G$ permutes the orbit. Each row of $A$ has $d$ entries $1$, so $1 + \varepsilon A$ is strictly diagonally dominant for $|\varepsilon| < 1/d$ and hence positive definite. It is not $(v - c)1 + cU$: its off-diagonal entries take the two values $\varepsilon$ and $0$, both occurring since there is a second orbit. A centred Gaussian measure on $\mathbb{R}^{N}$ with covariance $1 + \varepsilon A$ is invariant under the coordinate permutations of $G$, and its coordinates are observables $O_{i}$ permuted by $G$. Conversely (i). The symmetric $G$-invariant matrices are spanned by $1$ and the adjacency matrices of the orbits on unordered pairs, which gives the count. The check exhibits the ring of six cells with nearest-neighbour $A$, $\varepsilon = \tfrac14$: variances per unit norm $\tfrac54, \tfrac34, \tfrac12$ on the three cosine modes.

(iv) The transposition of planes $i$ and $j$ is an automorphism exchanging $a_{i}$ and $a_{j}$; invariance gives $E_{ij} = E_{ji}$, and $E$ is Hermitian for the conjugation of its scalars, so $E_{ij}$ is real; double transitivity gives $E_{ii} = n$ and $E_{ij} = e$ for $i \ne j$. For the pairing, $\omega(a_{i}a_{j}) = \omega(a_{j}a_{i})$ by the same transposition and $a_{i}a_{j} = -a_{j}a_{i}$ (the vectors $w_{i}$ anticommute, Theorem \claimref{tower.ladder-pairs-satisfy-car-vectors}), so it vanishes for $i \ne j$; and $a_{i}^{2} = Q(w_{i}) = 0$. From $0 \le E \le 1$ and the eigenvalues $n - e$ and $n + (N - 1)e$ of $E$: $0 \le n + (N - 1)e \le 1$, which is the stated bound, and $0 \le n - e \le 1$, which is the other. For quasi-free $\omega$ the pairing vanishes, so the state is gauge-invariant and $\mathrm{Var}(N_{f}) = \Tr(fEf(1 - E))$ (Proposition \claimref{synthesis.quasi-free-deposit-is-vacuum-plus-shot-noise}(i)); with $E = \alpha 1 + eU$ and $\Tr(f^{2}U) = \Tr(fUf) = \mathbf{1}^{T}f^{2}\mathbf{1}$, $\Tr(fUfU) = (\mathbf{1}^{T}f\mathbf{1})^{2}$ the trace is the displayed expression; for diagonal $f$, $\alpha(1 - \alpha) + e(1 - 2\alpha) = n(1 - n) + e^{2}$. The check realises four modes in $\Cl(4,4)$ and the state of one quantum in the symmetric mode, $\psi = (\sum_{j}v_{j})\prod_{j}(1 - n_{j})$, a pure quasi-free state: the transposition and the four-cycle of planes fix $\psi$; $E_{ij} = \tfrac14$ for all $i, j$, the upper end $e = (1 - n)/(N - 1)$; the pairing vanishes; the number covariances are $v = \tfrac{3}{16}$, $c = -\tfrac{1}{16}$; and the displayed variance agrees with the direct one for diagonal and for non-diagonal $f$. For the ladder states, permutation symmetry and the moments $\omega_{k}(n_{u_{i}}) = k/N$, $\omega_{k}(n_{u_{i}}n_{u_{j}}) = k(k - 1)/(N(N - 1))$ are stored (Theorem \claimref{synthesis.ladder-states-are-locally-product-with-non-gaussian-fluctuations}(ii)); since $n_{u_{i}}^{2} = n_{u_{i}}$, $v = \tfrac kN(1 - \tfrac kN)$, and the two combinations follow by arithmetic. The check recomputes them on two and three blocks.

(v) Given distinct $x, y$ and distinct $x', y'$ put $d = y - x$, $d' = y' - x'$, $\lambda = (Q(d')/Q(d))^{1/2}$ and $u = \lambda d - d'$. If $u = 0$ the map $p \mapsto \lambda(p - x) + x'$ carries $(x, y)$ to $(x', y')$. Otherwise $Q(\lambda d) = Q(d')$ gives $u \cdot (\lambda d + d') = 0$, so the reflection $z \mapsto -uzu^{-1}$ exchanges $\lambda d$ and $d'$, and $p \mapsto -u\,\lambda(p - x)\,u^{-1} + x'$ does it. So the similarities are transitive on ordered pairs of distinct points. An invariant kernel then takes one value. For weight $\Delta$, $K(x, y)\,Q(x - y)^{\Delta}$ is invariant, since $Q(gx - gy) = \lambda_{g}^{2}Q(x - y)$, hence constant. The check carries three rational pairs to three prescribed pairs with $\lambda = \tfrac73, \tfrac37, \tfrac{11}{9}$ and verifies the scaling of $Q(x - y)^{-\Delta}$ for $\Delta = 0, 1, -1, 2$. The index of a slice correlator homogeneous of degree $-2\Delta$ is $2\Delta$, and on the frozen phase the leading weight gives a non-negative index (Proposition \claimref{synthesis.frozen-phase-spectrum-not-red}); that proposition is cited for the reading and is not used in the proof of (v).

The consequence is (i) to (iii) read for a deposit whose units are placed in a slice: the state supplies $v - c$ and $c$, and every dependence on the wavenumber is in the quadratic forms $\sum_{i}f_{i}^{2}$ and $(\sum_{i}f_{i})^{2}$ of the placement, the second of which vanishes at nonzero $k$ for tilings and lattices and is the placement's structure factor otherwise. The reading that follows it is argued: it lists the two places a shape can come from, the state's invariance group and the placement, and does not exclude either.
